\section{Experiments}
\label{sec:experiments}

\subsection{Instances and Experimental Settings}
\label{subsec:settings}

% Describe datasets/instance generation, train-validation-test separation,
% hardware, software, random seeds, repetitions, and implementation details.
% Include a parameter table and justify important parameter choices.

\subsection{Baselines and Evaluation Metrics}
\label{subsec:baselines_metrics}

% Select strong and recent baselines from relevant categories. Explain how each
% was tuned and how fairness was maintained. Define every metric mathematically.

\subsection{Parameter and Ablation Studies}
\label{subsec:ablation}

% Isolate contributions from graph heterogeneity, relation/meta-path design,
% HGNN layers, action masking, reward terms, and PPO-specific mechanisms.

\subsection{Comparative Results}
\label{subsec:comparative_results}

% Report central tendency and variability across repeated runs. Use statistical
% tests where appropriate. Discuss both wins and failures without overstating them.

% Example modular table usage:
% \begin{table}[t]
  \centering
  \caption{Replace with a self-contained description of the reported results.}
  \label{tab:example_results}
  \begin{tabular}{lcc}
    \toprule
    Method & Metric 1 $\downarrow$ & Metric 2 $\uparrow$ \\
    \midrule
    Baseline & -- & -- \\
    Proposed method & -- & -- \\
    \bottomrule
  \end{tabular}
\end{table}


\subsection{Generalization and Runtime Analysis}
\label{subsec:generalization_runtime}

% Evaluate unseen sizes, distributions, dynamic conditions, or system layouts.
% Report decision latency and scaling behavior relevant to real-time deployment.
